\documentclass{article}
\usepackage{amsmath}
\usepackage{booktabs}
\usepackage{geometry}
\geometry{margin=2.5cm}

\title{Justification for Hidden Mass Energy Terms}
\date{}

\begin{document}
\maketitle

\section*{Justification for $T_{ma}$, $V_\text{spring}$, $\mathcal{D}_\text{damper}$}

\subsection*{Kinetic Energy $T_{ma}$}

Starting from Garc\'ia's own kinetic energy for $m_h$, which follows directly
from equation~(3) after applying the small-angle simplification
$\cos\Theta \approx 1$, $\sin\Theta \approx 0$ (Section~2.3):
%
\begin{equation}
    T_{mh} = \frac{1}{2}m_h\left[
        \left(\dot{X} - Y\dot{\Theta}\right)^2
        + \left(\dot{Y} - d\dot{\Theta}\right)^2
    \right] + \frac{1}{2}J_h\dot{\Theta}^2.
\end{equation}
%
The mass $m_h$ sits at position $Y$ along the crossarm, offset $d$
perpendicular to it. Its absolute velocities in the inertial frame are:
%
\begin{align}
    \dot{x}_{m_h} &= \dot{X} - Y\dot{\Theta}, \\
    \dot{y}_{m_h} &= \dot{Y} - d\dot{\Theta}.
\end{align}
%
The hidden mass $m_a$ follows the same geometric structure as $m_h$. It sits
at absolute position $Y + \delta_a$ along the crossarm, where $\delta_a$ is
the \emph{relative} displacement of $m_a$ from $m_h$. It carries the same
$d$-offset perpendicular to the crossarm, and moves with $m_h$ plus an
additional relative velocity $\dot{\delta}_a$:
%
\begin{table}[h]
    \centering
    \begin{tabular}{lcc}
        \toprule
         & $m_h$ & $m_a$ \\
        \midrule
        Position along crossarm  & $Y$                             & $Y + \delta_a$                         \\
        Absolute x-velocity      & $\dot{X} - Y\dot\Theta$         & $\dot{X} - (Y+\delta_a)\dot\Theta$     \\
        Absolute y-velocity      & $\dot{Y} - d\dot\Theta$         & $\dot{Y} + \dot\delta_a - d\dot\Theta$ \\
        Rotary inertia           & $J_h$ (physical body)           & none (point mass)                      \\
        \bottomrule
    \end{tabular}
    \caption{Kinematic comparison of $m_h$ and $m_a$ in the inertial frame.
             The substitutions $Y \to Y + \delta_a$ and
             $\dot{Y} \to \dot{Y} + \dot\delta_a$ are the only differences.
             $m_a$ is modelled as a point mass and therefore carries no
             rotary inertia term.}
\end{table}
%
The kinetic energy of $m_a$ therefore follows directly:
%
\begin{equation}
    T_{ma} = \frac{1}{2}m_a\left[
        \left(\dot{X} - (Y+\delta_a)\dot{\Theta}\right)^2
        + \left(\dot{Y} + \dot{\delta}_a - d\dot{\Theta}\right)^2
    \right].
\end{equation}
%
This is identical in form to $T_{mh}$, with the substitutions
$Y \to Y + \delta_a$, $\dot{Y} \to \dot{Y} + \dot{\delta}_a$,
and no rotary inertia term since $m_a$ is a point mass.

\subsection*{Potential Energy $V_\text{spring}$}

The spring of stiffness $k_a$ connects $m_a$ to $m_h$. It stores elastic
energy proportional to the relative linear displacement $\delta_a$ between
the two masses. By Hooke's law for a linear spring:
%
\begin{equation}
    V_\text{spring} = \frac{1}{2}k_a\,\delta_a^2.
\end{equation}
%
Since $\delta_a$ is defined as the relative displacement of $m_a$ from $m_h$,
no additional projection is needed.

\subsection*{Rayleigh Dissipation $\mathcal{D}_\text{damper}$}

The damper of coefficient $c_a$ connects $m_a$ to $m_h$. Following
Garc\'ia's use of the Rayleigh dissipation formalism (equation~(5)), the
dissipation function for a linear viscous damper is:
%
\begin{equation}
    \mathcal{D}_\text{damper} = \frac{1}{2}c_a\,\dot{\delta}_a^2,
\end{equation}
%
where $\dot{\delta}_a$ is the relative linear velocity across the damper.
This is directly analogous to Garc\'ia's $\frac{1}{2}c_y\dot{Y}^2$ term ---
both are linear viscous dampers with a linear relative velocity coordinate.
The Rayleigh dissipation function enters the Euler--Lagrange equations as
$\partial\mathcal{D}/\partial\dot{q}_j$, not as a potential energy.

\subsection*{Extended Energy Expressions}

The three extended energy expressions follow by adding the new terms
to Garc\'ia's verified baseline:
%
\begin{align}
    T                   &= T_\text{garcia} + T_{ma},                   \\
    V                   &= V_\text{garcia} + V_\text{spring},           \\
    \mathcal{D}         &= \mathcal{D}_\text{garcia} + \mathcal{D}_\text{damper}.
\end{align}
%
The equations of motion for the extended 4-DOF system with generalised
coordinates $q = [X,\,\Theta,\,Y,\,\delta_a]^\top$ are obtained by
applying the Euler--Lagrange equations:
%
\begin{equation}
    \frac{d}{dt}\frac{\partial T}{\partial \dot{q}_j}
    - \frac{\partial T}{\partial q_j}
    + \frac{\partial V}{\partial q_j}
    + \frac{\partial \mathcal{D}}{\partial \dot{q}_j}
    = f_j, \qquad q_j \in \{X,\,\Theta,\,Y,\,\delta_a\},
\end{equation}
%
where $f = [F_1+F_2,\;(F_1-F_2)L_b/2,\;F_Y,\;0]^\top$. The zero in the
last entry reflects that no external force acts directly on the hidden
mass $m_a$.

\end{document}